\begin{eabstract}
Low-Rank Adaptation (LoRA) is a practical method for adapting large pretrained models because it
freezes the base model and learns compact low-rank updates. Its simplicity is also a limitation:
the same rank and target modules are normally selected without regard to the size, distribution,
or complexity of the downstream data. AdaLoRA redistributes a parameter budget during training,
but introduces an allocation schedule, importance estimation, and sensitivity to short
low-resource runs.

This thesis presents \starlora, a dataset-aware extension of AdaLoRA for low-resource sequence
classification. Before training, Star-LoRA profiles sample count, class imbalance, normalized
label entropy, token sparsity, token-frequency skew, and input-embedding variance. A bounded policy
uses these statistics to select target rank, initial rank, regularization, and adapter target
modules. During training, a stability monitor records exponential-moving-average and
variance-normalized importance. The current prototype uses this score for analysis rather than
replacing AdaLoRA's internal allocator.

Experiments use Qwen2.5-0.5B on SST-2 subsets containing 100 to 10,000 training examples. In the
primary three-seed sweep, Star-LoRA selects average target ranks from 4.67 to 10 as the sample
budget grows. Relative to standard AdaLoRA, its largest accuracy improvement is 4.97 percentage
points at 5,000 examples. It reaches 91.55\% accuracy in this condition, compared with AdaLoRA's
86.58\%. Fixed-rank LoRA remains a strong and substantially faster baseline, showing that
dataset-aware allocation does not guarantee universal superiority.

An exploratory extension asks a large language model to propose a layer-wise adapter layout.
Claude-generated layouts reduce runtime but do not improve accuracy. Gemini-labeled runs are
excluded from LLM performance claims because the API requests failed and used heuristic fallback.
The results show that dataset-aware configuration can improve standard AdaLoRA, while also
identifying compute-aware selection, controlled ablation, and stability-gated allocation as
necessary future work.
\end{eabstract}
